\documentclass[fontsize=12pt,BCOR=5mm,DIV=12,parskip=half,listof=totoc,
               paper=a4,toc=bibliography,pointlessnumbers]{scrreprt}
               
\usepackage[utf8]{inputenc}

%% LANGUAGE SETTINGS
%
%\usepackage[ngerman]{babel} 	        % german language
%\usepackage[german=quotes]{csquotes} 	% correct quoting using \enquote{}
\usepackage[english]{babel}            % english language
\usepackage{csquotes} 	                % correct quoting using \enquote{}

%%%%%%%%%%%%%
%% ZITIERSTIL
%%%%%%%%%%%%%
%
% @stud: Zitierstil in package biblatex unten wählen
%
% NUMERIC Style - e. g. [12]
% style=numeric 
%
% IEEE Style - numeric kind of style 
% style=ieee 
%
% ALPHABETIC Style - e. g. [AB12]
% style=alphabetic 
%
% HARVARD Style 
% style=apa 
%
% CHICAGO Style 
% style=authoryear
%
% Position des Zitats:
%
% autocite=inline 
%
% (!!) FOOTNOTE POSITION NOT RECOMMENDED IN MINT DOMAIN:
% autocite=footnote
%
\usepackage[backend=biber, 
            autocite=inline, 
            style=authoryear,
            citestyle=authoryear,
            bibstyle=authoryear,
            citetracker=true,
            pagetracker=true,
            maxcitenames=2]{biblatex}
            
% Adding compatibility for parencite with page numbers
\DeclareFieldFormat{postnote}{#1}
\DeclareFieldFormat{multipostnote}{#1}

\usepackage{makeidx}                  % allows index generation
\usepackage{listings}	              %Format Listings properly
\usepackage{lipsum}                   % Blindtext
\usepackage{graphicx}                 % use various graphics formats
\usepackage[german]{varioref}         % nicer references \vref
\usepackage{caption}	              % better Captions
\usepackage{booktabs}                 % nicer Tabs
\usepackage[hidelinks=true]{hyperref} % keine roten Markierungen bei Links
\usepackage{fnpct}                    % Correct superscripts 
\usepackage{calc}                     % Used for extra space below footsepline, in particular
\usepackage{array}
\usepackage{acronym}
\usepackage{algorithm}
\usepackage{algpseudocode}
\usepackage{setspace}
\usepackage{tocloft}
\usepackage[T1]{fontenc}

% Definitionen und Commands
\newcommand{\indextype}{numeric}
\newcommand{\abs}{\par\vskip 0.2cm\goodbreak\noindent}
\newcommand{\nl}{\par\noindent}
\newcommand{\mcl}[1]{\mathcal{#1}}
\newcommand{\nowrite}[1]{}
\newcommand{\NN}{{\mathbb N}}
\newcommand{\imagedir}{img}
\newcommand{\TitelDerArbeit}[1]{\def\DerTitelDerArbeit{#1}\hypersetup{pdftitle={#1}}}
\newcommand{\AutorDerArbeit}[1]{\def\DerAutorDerArbeit{#1}\hypersetup{pdfauthor={#1}}}
\newcommand{\Firma}[1]{\def\DerNameDerFirma{#1}}
\newcommand{\Kurs}[1]{\def\DieKursbezeichnung{#1}}
\newcommand{\Abteilung}[1]{\def\DerNameDerAbteilung{#1}}
\newcommand{\Studiengangsleiter}[1]{\def\DerStudiengangsleiter{#1}}
\newcommand{\WissBetreuer}[1]{\def\DerWissBetreuer{#1}}
\newcommand{\FirmenBetreuer}[1]{\def\DerFirmenBetreuer{#1}}
\newcommand{\Bearbeitungszeitraum}[1]{\def\DerBearbeitungszeitraum{#1}}
\newcommand{\Abgabedatum}[1]{\def\DasAbgabedatum{#1}}
\newcommand{\Matrikelnummer}[1]{\def\DieMatrikelnummer{#1}}
\newcommand{\Studienrichtung}[1]{\def\DieStudienrichtung{#1}}
\newcommand{\ArtDerArbeit}[1]{\def\DieArtDerArbeit{#1}}
\newcommand{\Literaturverzeichnis}{Literaturverzeichnis}

% Page Layout
\oddsidemargin=0mm
\evensidemargin=0mm
\textwidth=159mm
\topmargin=-18mm
\headsep=10mm
\textheight=251mm
\footheight=15mm

\makeindex

%%%%%%%%%%%%%%%%%%%%%%%%%%%%%%%%%%%
% LITERATURVERZEICHNIS
%
\addbibresource{bibliography.bib}
\DefineBibliographyStrings{english}{andothers = {{et\,al\adddot}},}

% !TEX root =  master.tex

%%%%%%%%%%%%%%%%%%%%%%%%%%%%%%%%%%%%%%%%%%%%%%%%%%%%%%%%%%%%%%%%%%
%	ANLEITUNG: 
% Passen Sie gegebenenfalls alle Stellen im Dokument an, die mit 
% @stud 
% markiert sind.
%%%%%%%%%%%%%%%%%%%%%%%%%%%%%%%%%%%%%%%%%%%%%%%%%%%%%%%%%%%%%%%%%%

%%
%% @stud
%%
%% LANGUAGE SETTINGS
%\usepackage[ngerman]{babel} 	        % german language
%\usepackage[german=quotes]{csquotes} 	% correct quoting using \enquote{}
\usepackage[english]{babel}          % english language
\usepackage{csquotes} 	              % correct quoting using \enquote{}

\usepackage{makeidx}                  % allows index generation
\usepackage{listings}	                %Format Listings properly
\usepackage{lipsum}                   % Blindtext
\usepackage{graphicx}                 % use various graphics formats
\usepackage[german]{varioref}         % nicer references \vref
\usepackage{caption}	                % better Captions
\usepackage{booktabs}                 % nicer Tabs
\usepackage[hidelinks=true]{hyperref} % keine roten Markierungen bei Links
\usepackage{fnpct}                    % Correct superscripts 
\usepackage{calc}                     % Used for extra space below footsepline, in particular
\usepackage{array}
\usepackage{acronym}
\usepackage{algorithm}
\usepackage{algpseudocode}
\usepackage{setspace}
\usepackage{tocloft}

%% Schriftarten- und Zeichenpakete
\usepackage[T1]{fontenc}
\usepackage[utf8]{inputenc}

%%
%% @stud
%%
%%	FONT SELECTION: Schriftarten und Schriftfamilie
%%%%%%%%%%%%%
%% SCHRIFTART
%%%%%%%%%%%%%
% 0) without decomment: normal font families 
% ...
% 1) Latin Modern 
%\usepackage{lmodern}        
% 2) Times 
%\usepackage{mathptmx}         
% 3) Helvetica
%\usepackage[scaled=.92]{helvet} 
%%%%%%%%%%%%%%%%%%
%%	SCHRIFTFAMILIE
%%%%%%%%%%%%%%%%%%
% ohne Serifen
\renewcommand*{\familydefault}{\sfdefault}
\addtokomafont{disposition}{\sffamily}
%
% mit Serifen
%\renewcommand*{\familydefault}{\rmdefault}
%\addtokomafont{disposition}{\rmfamily}
%
% Typewriter
%\renewcommand*{\familydefault}{\ttdefault}
%\addtokomafont{disposition}{\ttfamily}

%%
%% @stud
%%
%% Uncomment the following lines to support hard URL breaks in bibliography 
%\apptocmd{\UrlBreaks}{\do\f\do\m}{}{}
%\setcounter{biburllcpenalty}{9000}% Kleinbuchstaben
%\setcounter{biburlucpenalty}{9000}% Großbuchstaben

%%
%% @stud
%%
%% FOOTNOTES: Count footnotes over chapters
%% \counterwithout{footnote}{chapter}

%	ACRONYMS
\makeatletter
\@ifpackagelater{acronym}{2015/03/20}
{\renewcommand*{\aclabelfont}[1]{\textbf{{\acsfont{#1}}}}}{}
\makeatother

%	LISTINGS
% @stud: ggf. Namen/Text anpassen (englisch)
\renewcommand{\lstlistingname}{Source Code} 
\renewcommand{\lstlistlistingname}{Source Code Listings}
\lstset{numbers=left,
	numberstyle=\tiny,
	captionpos=b,
	basicstyle=\ttfamily\small}

%	ALGORITHMS
% @stud: ggf. Namen/Text anpassen (englisch)
\renewcommand{\listalgorithmname}{Algorithms}
\floatname{algorithm}{Algorithms}

%	PAGE HEADER / FOOTER
%	Warning: There are some redefinitions throughout the master.tex-file!  DON'T CHANGE THESE REDEFINITIONS!
\RequirePackage[automark]{scrlayer-scrpage}
%alternatively with separation lines: \RequirePackage[automark,headsepline,footsepline]{scrlayer-scrpage}

\renewcommand{\chaptermarkformat}{}
\RedeclareSectionCommand[beforeskip=0pt]{chapter}
\clearpairofpagestyles

\ofoot[\rule{0pt}{\ht\strutbox+\dp\strutbox}\pagemark]{\rule{0pt}{\ht\strutbox+\dp\strutbox}\pagemark}
\ohead{\headmark}

\newcommand{\settingBibFootnoteCite}{
	\setlength{\bibparsep}{\parskip}		  % Add some space between biblatex entries in the bibliography
	\addbibresource{bibliography.bib}	    % Add file bibliography.bib as biblatex resource
	\DefineBibliographyStrings{ngerman}{andothers = {{et\,al\adddot}},}
}

\newcommand{\setTitlepage}{
	% !TEX root =  master.tex
% @stud: ggf. Namen/Text anpassen (englisch)
\begin{titlepage}
\begin{minipage}{\textwidth}
		\vspace{-1cm}
		\noindent \includegraphics[scale=0.450]{\imagedir/mercedes.png} \hfill \includegraphics[scale=0.3]{\imagedir/dhbw.png}
\end{minipage}
\vspace{1em}
%\sffamily
\begin{center}
	{\textsf{\large Duale Hochschule Baden-W\"urttemberg Stuttgart}}\\[4em]
	{\textsf{\textbf{\large{\DieArtDerArbeit} Thesis}}}\\[6mm]
	{\textsf{\textbf{\Large{}\DerTitelDerArbeit}}} \\[1.5cm]
	{\textsf{\textbf{\large{}Computer Science}}\\[6mm]
	\textsf{\textbf{\DieStudienrichtung}}}\vspace{10em}
	
	\begin{minipage}{\textwidth}
		\begin{tabbing}
		Wissenschaftlicher Betreuer: \hspace{0.85cm}\=\kill
		Author: \> \DerAutorDerArbeit \\[1.5mm]
		Student number: \> \DieMatrikelnummer \\[1.5mm]
		Company: \> \DerNameDerFirma  \\[1.5mm]
		Class: \> \DieKursbezeichnung \\[1.5mm]
		Course Administrator: \> \DerStudiengangsleiter \\[1.5mm]
		Academic Supervisor: \> \DerWissBetreuer \\[1.5mm]
		Technical Supervisor: \> \DerFirmenBetreuer \\[1.5mm]
		% Bearbeitungszeitraum: \> \DerBearbeitungszeitraum\\[1.5mm]
%		alternativ:\\[1.5mm]
		Submitted: \> \DasAbgabedatum	
		\end{tabbing}
	\end{minipage}
\end{center}
\end{titlepage}
	\pagenumbering{roman} % Römische Seitennummerierung
	\normalfont	
}

\newcommand{\initializeText}{
	\clearpage
	\ihead{\chaptername~\thechapter} % Neue Header-Definition
	\pagenumbering{arabic}           % Arabische Seitenzahlen
}

\newcommand{\initializeBibliography}{
	\ihead{}
	\printbibliography[title=\Litarature, heading=bibintoc] 
	\cleardoublepage
}

\newcommand{\initializeAppendix}{
	\appendix
  	\ihead{}
  	\cftaddtitleline{toc}{chapter}{Appendix}{}
}



\ArtDerArbeit{Bachelor} %-thesis
\TitelDerArbeit{Development of a PoC for a PCIe based Automotive Zonal Architecture}
\AutorDerArbeit{Philipp Gehrig}
\Abteilung{}
\Firma{Mercedes-Benz Group AG}
\Kurs{Informatik-ITA22}
\Studienrichtung{IT-Automotive}
\Matrikelnummer{5622763}
\Studiengangsleiter{Prof. Dr. Zoltán Zomotor}
\WissBetreuer{Klaus Wieland}
\FirmenBetreuer{Felix Bauer M.Sc.}
\Bearbeitungszeitraum{24.03.2025 - 09.09.2025}
\Abgabedatum{09. September 2025}

\begin{document}

\setTitlepage

%%%%%%%%%%%%%%%%%%%%%%%%%%%%%%%%%%%
% EHRENWÖRTLICHE ERKLÄRUNG
%
\clearpage
\chapter*{Declatation of Academic Integrity}

I hereby declare that I have written the present thesis titled \textit{\DerTitelDerArbeit} independently 
and have used no sources or aids other than those stated. Furthermore, I confirm that 
the submitted electronic version corresponds to the printed version.

\vspace{3cm}
Place, Date \hfill \DerAutorDerArbeit


\chapter*{Ehrenwörtliche Erklärung}

Hiermit erkläre ich, dass ich die vorliegende Arbeit mit dem Titel \textit{\DerTitelDerArbeit} 
selbstständig verfasst und keine anderen als die angegebenen Quellen und Hilfsmittel benutzt habe. 
Ich versichere zudem, dass die eingereichte elektronische Fassung mit der gedruckten Version übereinstimmt.

\vspace{3cm}
Ort, Datum \hfill \DerAutorDerArbeit 
\cleardoublepage  
%%%%%%%%%%%%%%%%%%%%%%%%%%%%%%%%%%%

%%%%%%%%%%%%%%%%%%%%%%%%%%%%%%%%%%%
%	KURZFASSUNG
%
\chapter*{Achkowledgement}

This chapter contains the Achkowledgement


\cleardoublepage 
%%%%%%%%%%%%%%%%%%%%%%%%%%%%%%%%%%%

%%%%%%%%%%%%%%%%%%%%%%%%%%%%%%%%%%%
% VERZEICHNISSE und ABSTRACT
%
\tableofcontents
\cleardoublepage

% Abbildungsverzeichnis
\phantomsection
\addcontentsline{toc}{chapter}{\listfigurename}
\listoffigures
\cleardoublepage

% Tabellenverzeichnis
\phantomsection
\addcontentsline{toc}{chapter}{\listtablename}
\listoftables
\cleardoublepage

% Listingsverzeichnis / Quelltextverzeichnis
%\lstlistoflistings
%\cleardoublepage

% Algorithmenverzeichnis
%\listofalgorithms
%\cleardoublepage

% Abkürzungsverzeichnis
% !TEX root =  master.tex
\clearpage
\chapter*{Acronyms}	
\addcontentsline{toc}{chapter}{Acronyms}

\begin{acronym}[YTMMM]
\setlength{\itemsep}{-\parsep}

%A
%B
%C
%D
%E
%F
%G
%H
%I
%J
%K
%L
%M
%N
%O
%P
%Q
%R
%S
%T
%U
%V
%W
%X
%Y
%Z


\end{acronym} 
\cleardoublepage

\onehalfspacing

% Kurzfassung / Abstract
\chapter*{Abstract} 
\cleardoublepage

\initializeText

%%%%%%%%%%%%%%%%%%%%%%%%%%%%%%%%%%%%%%%%%%%%%%%%%%%%%%%%%%%%%%%%%%%%%%%%%%%%%%%%%%%%%%%%%%
% KAPITEL UND ANHÄNGE
%
%%%%%%%%%%%%%%%%%%%%%%%%%%%%%%%%%%%%%%%%%%%%%%%%%%%%%%%%%%%%%%%%%%%%%%%%%%%%%%%%%%%%%%%%%%

%%%%%%%%%%%%%%%%%%%%%%%%%%%%%%%%%%%
% KAPITEL
%
% Einleitung

% mehrere Grundlagen- und Forschungs-Kapitel
\chapter{Introduction}
\chapter{State of the Art}
\label{chap:state-of-the-art}

\section{Peripheral Component Interconnect Express}
\label{sec:pcie}

The following section provides an overview of the \ac{PCIe} protocol, its architecture, its history, 
and its role in modern computing systems. Additionally, it will discuss the advancements in PCIe technology 
and its impact on data transfer rates and system performance.

\subsection{History and Fundamentals}
\label{subsec:pcie_history}

\ac{PCIe} is a high-speed serial computer expansion bus standard designed to 
replace the older \ac{PCI} and \ac{PCI-X} standards. In order to understand 
the evolution of \ac{PCIe}, it is essential to first understand the principles 
of its predecessors, \ac{PCI} and \ac{PCI-X}.

The development of modern computer bus architectures has been driven by the 
growing demand for higher bandwidth, lower latency, and better power 
efficiency \parencite{budruk2003pci}. This progression 
has led to significant advancements in how peripherals communicate 
with computer systems, culminating in the development of \ac{PCIe} as a 
dominant interconnect technology.

\subsubsection{Evolution of Computer Bus Architectures}
\label{subsubsec:pcie_evolution_bus}

The history of computer expansion bus technology represents a 
continuous evolution toward greater speeds, wider data paths, 
and enhanced functionality. Early personal computers utilized 
the Industry Standard Architecture (ISA) bus, introduced with 
the IBM PC/AT in 1984 \parencite{messmer1995pc}. 
While functional, ISA operated at just 8MHz with an 8-bit 
or 16-bit data path, creating a significant bottleneck 
as processor and peripheral speeds increased.

In response to these limitations, several interim standards emerged, 
including \ac{VLB} and the \ac{EISA} \parencite[pp. 9-10]{abbott2003pci}.
However, these solutions were eventually superseded by \ac{PCI}, 
which offered a more scalable and processor-independent approach to peripheral connectivity.

\subsubsection{PCI and PCI-X}
\label{subsubsec:pcie_pci_pcix}

\ac{PCI} was introduced by Intel in 1992 as a replacement for the \ac{ISA} 
bus \parencite{anderson1999pci}. \ac{PCI} operated as a parallel 
bus architecture with a 33 MHz clock speed and 32-bit data width, providing 
a maximum theoretical bandwidth of 133 MB/s \parencite{shanley2000pcix}. 
This represented a significant improvement over previous standards, enabling more 
efficient communication between the \ac{CPU} and peripheral devices.

The design philosophy behind \ac{PCI} was revolutionary for its time, introducing 
several key innovations \parencite[pp. 11, 92-105, 189]{abbott2003pci}:

\begin{itemize}
    \item \textbf{Processor independence}: Unlike earlier buses that were tightly 
    coupled to the CPU's architecture, PCI was designed to be processor-agnostic.
    \item \textbf{Bus mastering}: PCI devices could take control of the bus and 
    initiate transactions without CPU intervention, significantly improving multitasking capabilities.
    \item \textbf{Plug and Play}: PCI introduced automated device configuration, 
    eliminating the need for manual jumper settings that plagued earlier expansion cards.
\end{itemize}

As computing demands increased, the \ac{PCI} standard evolved into \ac{PCI-X} in 
1998 \parencite{shanley2000pcix}. \ac{PCI-X} maintained backward compatibility 
while increasing the bus frequency to 66 MHz, 133 MHz, and eventually 533 MHz in 
\ac{PCI-X} 2.0, enabling theoretical bandwidths of up to 4.3 GB/s. The \ac{PCI-X} 
standard was developed specifically to address the growing bandwidth demands of 
server applications, particularly for high-performance network and storage 
interfaces \parencite{shanley2000pcix}.

However, both \ac{PCI} and \ac{PCI-X} faced inherent limitations due to their parallel bus architecture:

\begin{itemize}
    \item \textbf{Signal integrity issues}: As frequencies increased, maintaining signal 
    integrity across multiple parallel lines became increasingly challenging 
    \parencite{wilen2003understanding}. At higher clock rates, 
    timing skew and cross-talk between adjacent traces created significant signal integrity problems.
    
    \item \textbf{Shared bandwidth}: All devices on the bus shared the same bandwidth, 
    creating bottlenecks in data-intensive applications 
    \parencite{budruk2003pci}. This shared resource model meant 
    that a single bandwidth-hungry device could degrade performance for all 
    other components on the same bus.
    
    \item \textbf{Complex routing}: The parallel architecture required complex 
    routing on motherboards and made high-speed operation difficult beyond 
    short distances \parencite{budruk2003pci}. The large number 
    of signal traces created challenges for motherboard design and limited 
    practical bus lengths.
    
    \item \textbf{Limited scalability}: The bus topology placed fundamental 
    limits on clock speeds and the number of supported devices 
    \parencite{anderson1999pci}. Each additional device added 
    electrical load to the bus, constraining maximum achievable frequencies.
    
    \item \textbf{Voltage issues}: Maintaining consistent voltage levels 
    across multiple devices operating at higher speeds became increasingly 
    difficult \parencite[pp.~112-115]{abbott2003pci}.
\end{itemize}

These limitations, coupled with the rapidly increasing performance requirements of 
modern computing systems, necessitated a fundamentally new approach to peripheral 
connectivity, ultimately leading to the development of \ac{PCIe}.

\subsubsection{Creation of PCIe}
\label{subsubsec:pcie_creation}

The development of \ac{PCIe} began in the late 1990s under the code name 
"Arapahoe" as a collaborative effort between several major technology companies, 
including Intel, Dell, IBM, and HP. In 2001, the technology was transferred to the \ac{PCI-SIG}, 
the industry organization responsible for specifying the \ac{PCI} 
family of standards. The first draft specification was completed in 2002, 
with the official \ac{PCIe} 1.0 specification released in 2003 
\parencite{pcisig2003}. This marked a revolutionary shift in computer 
architecture, moving from the shared parallel bus paradigm to a serial, 
point-to-point topology.

The fundamental innovation of \ac{PCIe} was its adoption of serialized 
data transmission over differential pairs, which provided several 
critical advantages \parencite{johnson2010pcie}:

\begin{itemize}
    \item \textbf{Superior signal integrity}: Differential signaling offered much 
    better noise immunity and allowed for higher clock rates.
    
    \item \textbf{Point-to-point connectivity}: Each device had a dedicated link to 
    either a switch or the root complex, eliminating shared bandwidth bottlenecks.
    
    \item \textbf{Scalable performance}: The introduction of "lanes" allowed 
    bandwidth to scale linearly with the number of lane pairs allocated to a device.
    
    \item \textbf{Lower pin count}: Serial transmission drastically reduced the 
    number of physical pins required for connectors.
\end{itemize}

From its inception, \ac{PCIe} was designed for backward compatibility at the 
software level, allowing existing operating systems and drivers to work with 
\ac{PCIe} devices with minimal modification \parencite{budruk2003pci}. 
This design philosophy ensured rapid industry adoption despite the radical changes 
in physical implementation.

\subsection{Generational Evolution of PCIe}
\label{subsec:pcie_evolution}

Since its introduction, \ac{PCIe} has undergone several generational advancements, 
each doubling the bandwidth of its predecessor while maintaining backward 
compatibility \parencite{pcisig2019}. This evolutionary approach 
has ensured the longevity and widespread adoption of the standard across 
various computing platforms. Table \ref{tab:pcie_generations} summarizes 
the key characteristics of each \ac{PCIe} generation.

\begin{table}[ht]
\centering
\caption{\ac{PCIe} Generations and Their Specifications}
\label{tab:pcie_generations}
\begin{tabular}{|l|c|c|c|c|}
\hline
\textbf{Generation} & \textbf{Year} & \textbf{Data Rate} & \textbf{Encoding} & \textbf{Bandwidth per Lane} \\
\hline
\acs{PCIe} 1.0 & 2003 & 2.5 \acs{GT/s} & 8b/10b & 250 MB/s \\ 
\acs{PCIe} 2.0 & 2007 & 5 \acs{GT/s} & 8b/10b & 500 MB/s \\ 
\acs{PCIe} 3.0 & 2010 & 8 \acs{GT/s} & 128b/130b & 985 MB/s \\ 
\acs{PCIe} 4.0 & 2017 & 16 \acs{GT/s} & 128b/130b & 1969 MB/s \\ 
\acs{PCIe} 5.0 & 2019 & 32 \acs{GT/s} & 128b/130b & 3938 MB/s \\ 
\acs{PCIe} 6.0 & 2022 & 64 \acs{GT/s} & \acs{PAM4}/\acs{FLIT} & 7877 MB/s \\ 
\acs{PCIe} 7.0* & 2025* & 128 \acs{GT/s} & \acs{PAM4}/\acs{FLIT} & 15754 MB/s \\
\hline
\end{tabular}
\begin{flushleft}
\small{* Projected specifications based on \ac{PCI-SIG} announcements \parencite{pcisig2023}}
\end{flushleft}
\end{table}

\ac{PCIe} 1.0 established the foundation with its 2.5 \ac{GT/s} data rate and 
8b/10b encoding scheme, resulting in 250 MB/s of effective bandwidth per direction 
per lane \parencite{budruk2003pci}. While modest by today's standards, 
this offered a significant improvement over \ac{PCI-X} for many applications, 
especially when aggregated across multiple lanes.

\ac{PCIe} 2.0, introduced in 2007, maintained the same encoding but doubled 
the transfer rate to 5 \ac{GT/s}, effectively doubling the bandwidth 
\parencite{pcisig2007}. This generation also introduced improved 
power management features, reflecting the growing importance of energy 
efficiency in computing systems \parencite{johnson2010pcie}.

A significant advancement came with \ac{PCIe} 3.0 in 2010, which not only 
increased the transfer rate to 8 \ac{GT/s} but also implemented a more 
efficient 128b/130b encoding scheme \parencite{pcisig2010}. 
This change reduced the encoding overhead from 20\% to approximately 
1.5\%, resulting in nearly 1 GB/s of bandwidth per lane in each direction. 
The introduction of this more efficient encoding was a critical innovation 
that significantly boosted effective throughput.

After a longer development cycle, \ac{PCIe} 4.0 was released in 2017, 
continuing the trend of doubling transfer rates to reach 16 \ac{GT/s} while 
maintaining the 128b/130b encoding, delivering approximately 2 GB/s per 
lane \parencite{pcisig2017}. This generation was particularly 
important for high-bandwidth applications such as graphics processing, 
machine learning acceleration, and high-performance storage.

\ac{PCIe} 5.0 followed relatively quickly in 2019, with transfer rates 
reaching 32 \ac{GT/s} with the same encoding scheme, effectively 
providing almost 4 GB/s per lane \parencite{pcisig2019}. 
The rapid progression from \ac{PCIe} 4.0 to 5.0—taking only about 
two years—reflected the increasing demands of data-intensive applications 
and the maturation of the development process.

\ac{PCIe} 6.0, finalized in 2022, represents a more fundamental 
shift in the standard's approach \parencite{pcisig2022}. 
While maintaining the doubled transfer rate pattern (64 \ac{GT/s}), 
it introduced \ac{PAM4} signaling and \ac{FLIT}-based encoding. \ac{PAM4} 
encodes two bits per symbol instead of one, allowing twice the data rate 
without increasing the fundamental frequency, which helps address signal 
integrity challenges at higher speeds. 
The \ac{FLIT} encoding scheme further optimizes transmission efficiency 
and incorporates \ac{FEC} to maintain reliability at these higher speeds.

Looking ahead, \ac{PCIe} 7.0 is currently under development with a projected 
release around 2025, aiming to continue the trend of doubling bandwidth to 
128 \ac{GT/s} using \ac{PAM4} signaling, potentially delivering approximately 
15.8 GB/s per lane in each direction \parencite{pcisig2023}. These 
continued advancements demonstrate \ac{PCIe}'s adaptability and its critical 
role in addressing the ever-increasing bandwidth demands of modern computing 
systems \parencite{mangla2019pcie}.

\subsection{Architectural Foundations}
\label{subsec:pcie_architecture_foundations}

The fundamental architecture of \ac{PCIe} consists of a layered protocol stack 
that separates different aspects of the communication process, providing 
flexibility and extensibility \parencite[pp.~22-26]{pcisig2003}. This layered 
approach consists of three primary protocol layers:

\begin{itemize}
    \item \textbf{Transaction Layer}: The uppermost layer responsible for generating 
    and processing \ac{TLP}s \parencite{budruk2003pci}. It manages end-to-end 
    communications, credit-based flow control, and quality of service mechanisms.
    \item \textbf{Data Link Layer}: The middle layer ensures reliable data transfer 
    by adding sequence numbers and \ac{CRC} protection to \ac{TLP}s, forming 
    \ac{DLLP}s \parencite{wilen2003understanding}. It handles packet 
    acknowledgment, retransmission of failed packets, and link power management.
    \item \textbf{Physical Layer}: The lowest layer deals with the electrical 
    aspects of transmitting and receiving serialized data \parencite{budruk2003pci}. 
    It manages lane initialization, training, and power management state transitions.
\end{itemize}

Unlike the shared bus architecture of \ac{PCI}, \ac{PCIe} implements 
a point-to-point topology built around three primary components: endpoints, 
switches, and root complexes \parencite[pp.~82-88]{solari2005pci}. 
This topology creates a hierarchical structure that enables more 
efficient data routing and resource allocation.

The root complex serves as the central connection point between 
the \ac{CPU}, memory subsystem, and \ac{PCIe} fabric, initiating 
configuration cycles and managing the enumeration process that 
identifies connected devices during system initialization 
\parencite{budruk2003pci}. \ac{PCIe} switches 
function as packet routing devices, directing traffic between 
multiple endpoints and the root complex \parencite[pp.~75-83]{wilen2003understanding}.

The addressing scheme in \ac{PCIe} uses a three-component system 
consisting of a bus number, device number, and function number 
(\ac{BDF}), similar to \ac{PCI} but expanded to accommodate the 
hierarchical topology \parencite[pp.~145-150]{solari2005pci}. 
This addressing scheme ensures that each function within the 
\ac{PCIe} hierarchy has a unique identifier, enabling precise 
routing of transactions throughout the fabric.
\section{Zonal Automotive Architecture}
\label{sec:automotive-architecture}
\chapter{Concept}
\label{chap:concept}

\section{Software Architecture}
\label{sec:architecture}
\chapter{Implementation}
\label{chap:implementation}

\section{Experimental Setup}
\label{sec:setup}
\section{Development of Software}
\label{sec:deployment}
\chapter{Evaluation}
\label{chap:evaluation}

\section{Benchmarking}
\label{sec:benchmarking}
\section{Reliability and Robustness}
\label{sec:reliability}



% Fazit und Ausblick
\chapter{Conclusion}
\chapter{Outlook}
\label{chap:outlook}

\section{Future Work}
\label{sec:future}
\section{Long Term Impact}
\label{sec:impact}
%%%%%%%%%%%%%%%%%%%%%%%%%%%%%%%%%%%

%%%%%%%%%%%%%%%%%%%%%%%%%%%%%%%%%%%
% ANHÄNGE
%
\initializeAppendix
\chapter{Appendix}
%%%%%%%%%%%%%%%%%%%%%%%%%%%%%%%%%%%

\singlespacing

\printbibliography 
\cleardoublepage

%%%%%%%%%%%%%%%%%%%%%%%%%%%%%%%%%%%
% INDEX
%
\addcontentsline{toc}{chapter}{Index}
\printindex

\end{document}
